\chapter{2013年辽宁卷（文）}

\section{单选题}

\begin{problem}
已知集合 \(A = \{0, 1, 2, 3, 4\}\)，\(B = \{x \mid |x| < 2\}\)，则 \(A \cap B =\)（\quad）

\choices
    {\( \{0\} \)}
    {\(\{0,1\}\)}
    {\(\{0,2\}\)}
    {\(\{0,1,2\}\)}
\end{problem}

\begin{answer}
B
\end{answer}

\begin{solution}
由 \(B=\{x\mid |x|<2\}=(-2,2)\)，而 \(A=\{0,1,2,3,4\}\)，所以
\(A\cap B=\{0,1\}\)，故选 B.
\end{solution}

\begin{problem}
复数的 \(z = \frac{1}{\i - 1}\) 模为（\quad）

\choices
    {\(\frac{1}{2}\)}
    {\(\frac{\sqrt{2}}{2}\)}
    {\(\sqrt{2}\)}
    {\(2\)}
\end{problem}

\begin{answer}
B
\end{answer}

\begin{solution}
\[
|z|=\left|\frac{1}{\i-1}\right|=\frac{1}{|\i-1|}
 =\frac{1}{\sqrt{(-1)^2+1^2}}=\frac{\sqrt{2}}{2}
\]
故选 B.
\end{solution}

\begin{problem}
已知点 \(A(1,3)\)，\(B(4,-1)\)，则与向量 \(\overrightarrow{AB}\) 同方向的单位向量为（\quad）

\choices
    {\(\left(\frac{3}{5}, -\frac{4}{5}\right)\)}
    {\( \left(\frac{4}{5},-\frac{3}{5}\right) \)}
    {\( \left(\frac{3}{5},\frac{4}{5}\right) \)}
    {\(\left(-\frac{3}{5}, \frac{4}{5}\right)\)}
\end{problem}

\begin{answer}
A
\end{answer}

\begin{solution}
\(\overrightarrow{AB}=(4-1,-1-3)=(3,-4)\)，且
\(\lvert\overrightarrow{AB}\rvert=\sqrt{3^2+(-4)^2}=5\). 因此与
\(\overrightarrow{AB}\) 同方向的单位向量为
\[
\frac{\overrightarrow{AB}}{\lvert\overrightarrow{AB}\rvert}
 =\left(\frac35,-\frac45\right)
\]
故选 A.
\end{solution}

\begin{problem}
下面是关于公差 \(d > 0\) 的等差数列 \(\{a_{n}\}\) 的四个命题：

\begin{circlelist}
\item \(p_{1}\)：数列 \(\{a_{n}\}\) 是递增数列；
\item \(p_{2}\)：数列 \(\{na_{n}\}\) 是递增数列；
\item \(p_{3}\)：数列 \(\left\{\frac{a_{n}}{n}\right\}\) 是递增数列；
\item \(p_{4}\)：数列 \(\{a_{n}+3nd\}\) 是递增数列.
\end{circlelist}

其中的真命题为（\quad）

\choices
    {\(p_1\)，\(p_2\)}
    {\(p_3\)，\(p_4\)}
    {\(p_2\)，\(p_3\)}
    {\(p_1\)，\(p_4\)}
\end{problem}

\begin{answer}
D
\end{answer}

\begin{solution}
设 \(a_n=a_1+(n-1)d\)，其中 \(d>0\). 则
\[
a_{n+1}-a_n=d>0
\]
所以 \(p_1\) 正确.

对 \(p_2\)，有
\[
(n+1)a_{n+1}-na_n=a_n+(n+1)d=a_1+2nd
\]
该式不一定为正，例如可取 \(a_1=-100\)，\(d=1\)，\(n=1\)，故 \(p_2\) 不一定正确.

对 \(p_3\)，有
\[
\frac{a_{n+1}}{n+1}-\frac{a_n}{n}
 =\frac{nd-a_n}{n(n+1)}
\]
该式不一定为正，例如可取 \(a_1=100\)，\(d=1\)，故 \(p_3\) 不一定正确.

最后
\[
(a_{n+1}+3(n+1)d)-(a_n+3nd)=4d>0
\]
所以 \(p_4\) 正确. 真命题为 \(p_1\)，\(p_4\)，故选 D.
\end{solution}

\begin{problem}
某班的全体学生参加英语测试，成绩的频率分布直方图如图，数据的分组依次为：\([20,40)\)，\([40,60)\)，\([60,80)\)，\([80,100]\). 若低于 \(60\) 分的人数是 \(15\)，则该班的学生人数是（\quad）

\bitmapfigure[max width=0.42\linewidth]{img/2013/liaoning_liberal/q5_fig1.png}

\choices
    {\(45\)}
    {\(50\)}
    {\(55\)}
    {\(60\)}
\end{problem}

\begin{answer}
B
\end{answer}

\begin{solution}
由直方图，\([20,40)\) 和 \([40,60)\) 两组的频率分别为
\(20\times0.005=0.1\) 与 \(20\times0.01=0.2\)，所以低于 \(60\) 分的频率为 \(0.3\). 设全班有 \(N\) 人，则
\[
0.3N=15
\]
解得 \(N=50\). 故选 B.
\end{solution}

\begin{problem}
在 \(\triangle ABC\) 中，内角 \(A\)，\(B\)，\(C\) 所对的边分别为 \(a\)，\(b\)，\(c\)，若 \(a \sin B \cos C + c \sin B \cos A = \frac{1}{2} b\)，且 \(a > b\)，则 \(\angle B =\)（\quad）

\choices
    {\(\frac{\pi}{6}\)}
    {\(\frac{\pi}{3}\)}
    {\(\frac{2\pi}{3}\)}
    {\(\frac{5\pi}{6}\)}
\end{problem}

\begin{answer}
A
\end{answer}

\begin{solution}
由正弦定理，存在正数 \(2R\)，使得
\(a=2R\sin A\)，\(b=2R\sin B\)，\(c=2R\sin C\). 因而
\[
a\cos C+c\cos A
 =2R(\sin A\cos C+\sin C\cos A)
 =2R\sin(A+C)=2R\sin B=b
\]
题设等式化为 \(b\sin B=\frac12b\)，所以
\(\sin B=\frac12\). 由于 \(0<B<\pi\)，故
\(B=\frac\pi6\) 或 \(B=\frac{5\pi}6\).

又因 \(a>b\)，由正弦定理得 \(\sin A>\sin B=\frac12\). 若
\(B=\frac{5\pi}6\)，则 \(A<\frac\pi6\)，从而
\(\sin A<\frac12\)，矛盾. 因此
\[
B=\frac\pi6
\]
故选 A.
\end{solution}

\begin{problem}
已知函数 \( f(x) = \ln (\sqrt{1 + 9x^2} - 3x) + 1 \)，则 \( f(\lg 2) + f\left(\lg \frac{1}{2}\right) = \)（\quad）

\choices
    {\(-1\)}
    {\(0\)}
    {\(1\)}
    {\(2\)}
\end{problem}

\begin{answer}
D
\end{answer}

\begin{solution}
令 \(t=\lg2\)，则 \(\lg\frac12=-t\). 于是
\[
f(t)+f(-t)=\ln\left[(\sqrt{1+9t^2}-3t)(\sqrt{1+9t^2}+3t)\right]+2=\ln1+2=2
\]
故选 D.
\end{solution}

\begin{problem}
执行如图所示的程序框图，若输入 \(n = 8\)，则输出 \(S =\)（\quad）

\bitmapfigure[max width=0.24\linewidth]{img/2013/liaoning_liberal/q8_fig1.png}

\choices
    {\(\frac{4}{9}\)}
    {\(\frac{6}{7}\)}
    {\(\frac{8}{9}\)}
    {\(\frac{10}{11}\)}
\end{problem}

\begin{answer}
A
\end{answer}

\begin{solution}
程序从 \(i=2\) 开始，每次循环把 \(i\) 增加 \(2\)，当 \(i\leqslant8\) 时累加. 因此
\[
S=\frac1{2^2-1}+\frac1{4^2-1}+\frac1{6^2-1}+\frac1{8^2-1}
\]
利用
\[
\frac1{i^2-1}=\frac12\left(\frac1{i-1}-\frac1{i+1}\right)
\]
得
\[
S=\frac12\left(1-\frac13+\frac13-\frac15+\frac15-\frac17+\frac17-\frac19\right)
 =\frac12\left(1-\frac19\right)=\frac49
\]
故选 A.
\end{solution}

\begin{problem}
已知点 \(O(0,0)\)，\(A(0,b)\)，\(B(a,a^3)\). 若 \(\triangle OAB\) 为直角三角形，则必有（\quad）

\choices
    {\(b = a^3\)}
    {\(b = a^3 + \frac{1}{a}\)}
    {\((b - a^3)\left(b - a^3 - \frac{1}{a}\right) = 0\)}
    {\( |b - a^3| + \left|b - a^3 - \frac{1}{a}\right| = 0 \)}
\end{problem}

\begin{answer}
C
\end{answer}

\begin{solution}
三角形非退化，所以直角不能在 \(O\) 点；只需分别讨论直角在 \(A\) 或 \(B\).

若 \(\angle A=90^\circ\)，则
\[
\overrightarrow{AO}=(0,-b)
\]
又 \(\overrightarrow{AB}=(a,a^3-b)\)，
从而
\[
\overrightarrow{AO}\cdot\overrightarrow{AB}=b(b-a^3)=0
\]
因三角形非退化，得 \(b=a^3\).

若 \(\angle B=90^\circ\)，则
\[
\overrightarrow{BO}=(-a,-a^3)
\]
又 \(\overrightarrow{BA}=(-a,b-a^3)\)，
所以
\[
\overrightarrow{BO}\cdot\overrightarrow{BA}
 =a^2-a^3(b-a^3)
 =-a^3\left(b-a^3-\frac1a\right)=0
\]
因三角形非退化，得 \(b=a^3+\frac1a\). 因此必有
\[
\left(b-a^3\right)\left(b-a^3-\frac1a\right)=0
\]
故选 C.
\end{solution}

\begin{problem}
已知直三棱柱 \(ABC-A_{1}B_{1}C_{1}\) 的 \(6\) 个顶点都在球 \(O\) 的球面上. 若 \(AB = 3\)，\(AC = 4\)，\(AB \perp AC\)，\(AA_{1} = 12\)，则球 \(O\) 的半径为（\quad）

\choices
    {\(\frac{3\sqrt{17}}{2}\)}
    {\(2\sqrt{10}\)}
    {\(\frac{13}{2}\)}
    {\(3\sqrt{10}\)}
\end{problem}

\begin{answer}
C
\end{answer}

\begin{solution}
底面 \(\triangle ABC\) 在 \(A\) 点直角，所以
\[
BC=\sqrt{AB^2+AC^2}=\sqrt{3^2+4^2}=5
\]
底面三角形的外接圆半径为 \(BC/2=5/2\). 直三棱柱的球心在上下两个底面外接圆圆心连线的中点处，因此球半径满足
\[
R^2=\left(\frac52\right)^2+\left(\frac{AA_1}{2}\right)^2
 =\frac{25}{4}+6^2=\frac{169}{4}
\]
故 \(R=\frac{13}{2}\)，选 C.
\end{solution}

\begin{problem}
已知椭圆 \(C:\frac{x^2}{a^2} + \frac{y^2}{b^2} = 1\)（\(a > b > 0\)）的左焦点为 \(F\)，\(C\) 与过原点的直线相交于 \(A\)，\(B\) 两点，连接 \(AF\)，\(BF\)，若 \(|AB| = 10\)，\(|BF| = 8\)，\(\cos \angle ABF = \frac{4}{5}\)，则 \(C\) 的离心率为（\quad）

\choices
    {\(\frac{3}{5}\)}
    {\( \frac{5}{7} \)}
    {\( \frac{4}{5} \)}
    {\( \frac{6}{7} \)}
\end{problem}

\begin{answer}
B
\end{answer}

\begin{solution}
由余弦定理，在 \(\triangle ABF\) 中
\[
AF^2=AB^2+BF^2-2\cdot AB\cdot BF\cos\angle ABF
 =10^2+8^2-2\cdot10\cdot8\cdot\frac45=36
\]
所以 \(AF=6\).

设椭圆的右焦点为 \(F'\). 因直线 \(AB\) 过原点，\(A\)，\(B\) 关于原点对称，所以 \(AF'=BF\). 由椭圆定义
\[
2a=AF+AF'=AF+BF=6+8=14
\]
从而 \(a=7\).

设 \(A=(x,y)\)，则 \(B=(-x,-y)\)，且 \(AB=10\) 给出
\(x^2+y^2=25\). 设左焦点 \(F=(-c,0)\)，由
\[
AF^2=(x+c)^2+y^2=x^2+y^2+c^2+2cx
\]
和 \(BF^2=(c-x)^2+y^2\)，得
\[
AF^2-BF^2=4cx=36-64=-28
\]
于是 \(2cx=-14\)，代入 \(AF^2=36\) 得
\[
36=25+c^2-14=c^2+11
\]
所以 \(c=5\). 因而离心率为
\[
e=\frac ca=\frac57
\]
故选 B.
\end{solution}

\begin{problem}
已知函数 \( f(x) = x^2 - 2(a + 2)x + a^2 \)，\( g(x) = -x^2 + 2(a - 2)x - a^2 + 8 \)，设 \( H_1(x) = \max \{f(x), g(x)\} \)，\( H_2(x) = \min \{f(x), g(x)\} \)（\(\max \{p, q\}\) 表示 \(p\)，\(q\) 中的较大值，\(\min \{p, q\}\) 表示 \(p\)，\(q\) 中的较小值）. 记 \( H_1(x) \) 的最小值为 \(A\)，\( H_2(x) \) 的最大值为 \(B\)，则 \( A - B = \)（\quad）

\choices
    {\(16\)}
    {\(-16\)}
    {\(a^2 - 2a - 16\)}
    {\(a^2 + 2a - 16\)}
\end{problem}

\begin{answer}
B
\end{answer}

\begin{solution}
两函数之差为
\[
f(x)-g(x)=2\bigl((x-a)^2-4\bigr)
\]
故交点横坐标为 \(a-2\) 与 \(a+2\)，且在区间 \([a-2,a+2]\) 内有 \(f\leqslant g\).

函数 \(H_1=\max\{f,g\}\) 在左侧取 \(f\)，中间取 \(g\)，右侧取 \(f\). 因为
\(f(x)=(x-a-2)^2-4a-4\)，其在右侧交点处取得最小值，故
\[
A=H_1(a+2)=f(a+2)=-4a-4
\]

函数 \(H_2=\min\{f,g\}\) 在中间取 \(f\)，两侧取 \(g\). 在中间区间上，\(f\)
在左端点取得最大值；左侧的 \(g\) 也在左端点取得最大值，而右侧的 \(g\) 在右端点取得最大值
\(g(a+2)=-4a-4<12-4a\)，因此
\[
B=H_2(a-2)=f(a-2)=g(a-2)=12-4a
\]
所以
\[
A-B=(-4a-4)-(12-4a)=-16
\]
故选 B.
\end{solution}

\section{填空题}

\begin{problem}
某几何体的三视图如图所示.

则该几何体的体积是\fillinblank{}.

\noindent
\leftskip=0pt\relax
\bitmapfigure[float=inline,width=0.40\linewidth]{img/2013/liaoning_liberal/q13_fig1.png}
\end{problem}

\begin{answer}
\(16\pi - 16\)
\end{answer}

\begin{solution}
由三视图可知，该几何体是高为 \(4\)、底面半径为 \(2\) 的圆柱去掉一个底面边长为 \(2\)、高为 \(4\) 的正四棱柱. 因此体积为
\[
V=\pi\cdot2^2\cdot4-2^2\cdot4=16\pi-16
\]
\end{solution}

\begin{problem}
已知等比数列 \(\{a_n\}\) 是递增数列，\(S_n\) 是 \(\{a_n\}\) 的前 \(n\) 项和. 若 \(a_1\)，\(a_3\) 是方程 \(x^2 - 5x + 4 = 0\) 的两个根，则 \(S_6 =\)\fillinblank{}.
\end{problem}

\begin{answer}
63
\end{answer}

\begin{solution}
方程 \(x^2-5x+4=0\) 的两根为 \(1,4\). 设等比数列首项为 \(a_1\)，公比为 \(q\).
若 \(a_1=1\)，\(a_3=4\)，则 \(q^2=4\). 由于数列递增且首项为正，必须有 \(q>1\)，故 \(q=2\).
若 \(a_1=4\)，\(a_3=1\)，则 \(q^2=1/4\)，不可能使首项为正的等比数列递增. 所以
\[
a_1=1
\]
从而
\[
S_6=\frac{1(2^6-1)}{2-1}=63
\]
\end{solution}

\begin{problem}
已知 \(F\) 为双曲线 \(C:\frac{x^2}{9} - \frac{y^2}{16} = 1\) 的左焦点，\(P\)，\(Q\) 为 \(C\) 上的点. 若 \(PQ\) 的长等于虚轴长的 \(2\) 倍，点 \(A(5,0)\) 在线段 \(PQ\) 上，则 \(\triangle PQF\) 的周长为\fillinblank{}.
\end{problem}

\begin{answer}
44
\end{answer}

\begin{solution}
双曲线的实半轴、虚半轴和半焦距分别为 \(a=3\)，\(b=4\)，\(c=5\)，右焦点为
\(F_2=(5,0)=A\)，而题设的左焦点为 \(F=(-5,0)\). 由题设
\[
PQ=2\times 2b=16
\]
点 \(A=F_2\) 在线段 \(PQ\) 上，所以
\[
PF_2+QF_2=PQ=16
\]
下面说明 \(P\)，\(Q\) 确在右支上. 若 \(PQ\) 不是竖直线，设其方程为
\(y=m(x-5)\)，并令 \(u=m^2\). 代入双曲线方程，令 \(t=x-5\)，得
\[
(16-9u)t^2+160t+256=0
\]
因 \(A\) 位于两交点之间，须有 \(u>\frac{16}{9}\). 设两交点的横坐标为
\(x_1\)，\(x_2\)，由求根公式得
\[
|x_1-x_2|=\frac{96\sqrt{1+u}}{9u-16}
\]
由于对应点的纵坐标差为 \(m(x_1-x_2)\)，所以
\[
PQ=|x_1-x_2|\sqrt{1+m^2}
=\frac{96(1+u)}{9u-16}
\]
结合 \(PQ=16\)，得 \(u=\frac{22}{3}\). 此时两交点的横坐标为
\[
x=\frac{33\pm8\sqrt3}{5}>3
\]
故均在右支上. 若 \(PQ\) 为竖直线，则两交点为
\((5,\pm\frac{16}{3})\)，弦长为 \(\frac{32}{3}\ne16\)，也不符合题设. 因而
双曲线定义给出
\[
PF-PF_2=QF-QF_2=2a=6
\]
因此
\[
PF+QF=(PF_2+6)+(QF_2+6)=16+12=28
\]
所求三角形周长为
\[
PQ+PF+QF=16+28=44
\]
\end{solution}

\begin{problem}
为了考察某校各班参加课外书法小组的人数，从全校随机抽取 \(5\) 个班级，把每个班级参加该小组的人数作为样本数据. 已知样本平均数为 \(7\)，样本方差为 \(4\)，且样本数据互不相同，则样本数据中的最大值为\fillinblank{}.
\end{problem}

\begin{answer}
10
\end{answer}

\begin{solution}
设样本数据为 \(x_1\)，\(\ldots\)，\(x_5\)，令 \(y_i=x_i-7\). 由平均数和方差得
\[
\sum_{i=1}^5y_i=0
\]
并且
\[
\sum_{i=1}^5y_i^2=5\times4=20
\]
设最大值为 \(7+k\)，则其余四个偏差之和为 \(-k\). 由柯西不等式，四个偏差的平方和至少为
\(k^2/4\)，所以
\[
20\geqslant k^2+\frac{k^2}{4}=\frac54k^2
\]
从而 \(k\leqslant4\). 若 \(k=4\)，上式取等号，四个其余偏差必须全为 \(-1\)，这与样本数据互不相同矛盾. 又因 \(x_i\) 表示人数，\(k\) 为整数，所以 \(k\leqslant3\)，最大值不超过 \(10\).

数据 \(4\)，\(6\)，\(7\)，\(8\)，\(10\) 的平均数为 \(7\)，且
\[
\frac{(4-7)^2+(6-7)^2+(7-7)^2+(8-7)^2+(10-7)^2}{5}=4
\]
满足所有条件，故最大值为 \(10\).
\end{solution}

\section{解答题}

\begin{problem}
\begin{enumerate}
\item 设向量 \(\bs{a} = (\sqrt{3}\sin x, \sin x)\)，\(\bs{b} = (\cos x, \sin x)\)，\(x \in \left[0, \frac{\pi}{2}\right]\). 若 \(\left|\bs{a}\right| = \left|\bs{b}\right|\)，求 \(x\) 的值；
\item 设函数 \(f(x) = \bs{a} \cdot \bs{b}\)，求 \(f(x)\) 的最大值.
\end{enumerate}
\end{problem}

\begin{solution}
\begin{enumerate}
\item
\[
|\bs{a}|^2=3\sin^2x+\sin^2x=4\sin^2x
\]
并且
\[
|\bs{b}|^2=\cos^2x+\sin^2x=1
\]
由 \(|\bs{a}|=|\bs{b}|\) 得 \(4\sin^2x=1\). 因
\(x\in[0,\frac\pi2]\)，所以 \(\sin x\geqslant0\)，从而
\[
x=\frac\pi6
\]

\item
\[
f(x)=\bs{a}\cdot\bs{b}=\sqrt3\sin x\cos x+\sin^2x
 =\frac12+\frac12(\sqrt3\sin2x-\cos2x)
\]
又
\[
\sqrt3\sin2x-\cos2x=2\sin\left(2x-\frac\pi6\right)
\]
当 \(x\in[0,\frac\pi2]\) 时，取 \(x=\frac\pi3\) 可使上式取得最大值 \(2\)，因此
\[
\max f(x)=\frac12+\frac12\times2=\frac32
\]
\end{enumerate}
\end{solution}

\begin{problem}
\begin{enumerate}
\item 如图，\(AB\) 是圆的直径，\(PA\) 垂直圆所在的平面，\(C\) 是圆上的点. 求证：\(BC \perp\) 平面 \(PAC\)；
\item 设 \(Q\) 为 \(PA\) 的中点，\(G\) 为 \(\triangle AOC\) 的重心，求证：\(QG\parallel\) 平面 \(PBC\).
\end{enumerate}
\bitmapfigure[float=inline,max width=0.20\linewidth]{img/2013/liaoning_liberal/q18_fig1.png}
\end{problem}

\begin{solution}
\begin{enumerate}
\item 由 \(AB\) 是圆的直径，得 \(\angle ACB=90^\circ\)，所以
\(BC\perp AC\). 又因为 \(PA\perp\) 平面 \(ABC\)，且 \(BC\) 在平面 \(ABC\) 内，所以
\(BC\perp PA\). 直线 \(PA\)，\(AC\) 是平面 \(PAC\) 内相交的两条直线，故
\[
BC\perp\text{ 平面 }PAC
\]

\item 以 \(A\) 为参考点，记
\(\overrightarrow{AP}=\bs{p}\)，\(\overrightarrow{AB}=\bs{b}\)，\(\overrightarrow{AC}=\bs{c}\). 因 \(O\)
是 \(AB\) 的中点，\(\overrightarrow{AO}=\frac12\bs{b}\). 又 \(Q\) 是 \(PA\) 的中点，\(G\) 是
\(\triangle AOC\) 的重心，所以
\[
\overrightarrow{AQ}=\frac12\bs{p}
\]
并且 \(\overrightarrow{AG}=\frac13\left(\frac12\bs{b}+\bs{c}\right)\)，
因此
\[
\overrightarrow{QG}=\frac16\bs{b}+\frac13\bs{c}-\frac12\bs{p}
 =\frac16(\bs{b}-\bs{p})+\frac13(\bs{c}-\bs{p})
 =\frac16\overrightarrow{PB}+\frac13\overrightarrow{PC}
\]
该向量是平面 \(PBC\) 内两个向量的线性组合，故
\[
QG\parallel\text{ 平面 }PBC
\]
\end{enumerate}
\end{solution}

\begin{problem}
现有 \(6\) 道题，其中 \(4\) 道甲类题，\(2\) 道乙类题，张同学从中任取 \(2\) 道题解答. 试求：

\begin{enumerate}
\item 所取的 \(2\) 道题都是甲类题的概率；
\item 所取的 \(2\) 道题不是同一类题的概率.
\end{enumerate}
\end{problem}

\begin{solution}
从 \(6\) 道题中任取 \(2\) 道，基本事件总数为
\(\mathrm C_6^2=15\).
\begin{enumerate}
\item 所取两题都是甲类题的基本事件数为 \(\mathrm C_4^2=6\)，所以概率为
\[
\frac{\mathrm C_4^2}{\mathrm C_6^2}=\frac6{15}=\frac25
\]

\item 所取两题不是同一类题，即一题甲类、一题乙类，基本事件数为
\(\mathrm C_4^1\mathrm C_2^1=8\)，所以概率为
\[
\frac{\mathrm C_4^1\mathrm C_2^1}{\mathrm C_6^2}=\frac8{15}
\]
\end{enumerate}
\end{solution}

\begin{problem}
如图，抛物线 \(C_{1}:x^{2} = 4y\)，\(C_{2}:x^{2} = -2py\)（\(p > 0\)）. 点 \(M(x_{0}, y_{0})\) 在抛物线 \(C_{2}\) 上，过 \(M\) 作 \(C_{1}\) 的切线，切点为 \(A\)，\(B\)（\(M\) 为原点 \(O\) 时，\(A\)，\(B\) 重合于 \(O\)）. 当 \(x_{0} = 1 - \sqrt{2}\) 时，切线 \(MA\) 的斜率为 \(-\frac{1}{2}\).

\begin{enumerate}
\item 求 \(p\) 的值；
\item 当 \(M\) 在 \(C_{2}\) 上运动时，求线段 \(AB\) 中点 \(N\) 的轨迹方程（\(A\)，\(B\) 重合于 \(O\) 时，中点为 \(O\)）.
\end{enumerate}

\FigureLayoutDeclare{img/2013/liaoning_liberal/q20_fig1.png}{5.8cm}{0bp 0bp 0bp 0bp}
\bitmapfigure[max width=0.28\linewidth]{img/2013/liaoning_liberal/q20_fig1.png}
\end{problem}

\begin{solution}
\begin{enumerate}
\item 将 \(C_1\) 写成 \(y=x^2/4\)，其在点 \((x_A,y_A)\) 处切线的斜率为
\(x_A/2\). 由斜率为 \(-1/2\)，得
\[
x_A=-1
\]
并且 \(y_A=\frac14\)，
切线 \(MA\) 为
\[
y-\frac14=-\frac12(x+1),
\qquad\text{即}\qquad y=-\frac12x-\frac14
\]
当 \(x_0=1-\sqrt2\) 时，因 \(M\) 在该切线上且在 \(C_2\) 上，故
\[
-\frac{(1-\sqrt2)^2}{2p}
 =-\frac12(1-\sqrt2)-\frac14
 =-\frac{3-2\sqrt2}{4}
\]
于是 \(p=2\).

\item 设切点 \(A\)，\(B\) 分别对应参数 \(m_1\)，\(m_2\)，即
\[
A=(2m_1,m_1^2)
\]
并且 \(B=(2m_2,m_2^2)\).
抛物线 \(C_1\) 在该点的切线方程为 \(y=mx-m^2\). 因切线经过
\(M(x_0,y_0)\)，所以 \(m_1\)，\(m_2\) 是方程
\[
m^2-x_0m+y_0=0
\]
的两根，从而
\[
m_1+m_2=x_0
\]
并且 \(m_1m_2=y_0\).
设 \(N=(x,y)\) 为 \(AB\) 中点，则
\[
x=m_1+m_2=x_0
\]
且
\[
y=\frac{m_1^2+m_2^2}{2}
 =\frac{(m_1+m_2)^2-2m_1m_2}{2}
 =\frac{x_0^2-2y_0}{2}
\]
由 \(p=2\)，曲线 \(C_2\) 为 \(x_0^2=-4y_0\)，所以
\[
y=\frac{x_0^2+x_0^2/2}{2}=\frac34x_0^2
\]
又 \(x=x_0\)，故轨迹方程为
\[
x^2=\frac43y
\]
当 \(M=O\) 时 \(x_0=y_0=0\)，上式也给出 \(N=O\)，故不需另加条件.
\end{enumerate}
\end{solution}

\begin{problem}
\begin{enumerate}
\item 证明：当 \(x \in [0,1]\) 时，\(\frac{\sqrt{2}}{2}x \leqslant \sin x \leqslant x\)；
\item 若不等式 \(ax + x^2 + \frac{x^3}{2} + 2(x + 2)\cos x \leqslant 4\) 对 \(x \in [0,1]\) 恒成立，求实数 \(a\) 的取值范围.
\end{enumerate}
\end{problem}

\begin{solution}
\begin{enumerate}
\item 设 \(h(x)=x-\sin x\). 在 \([0,1]\) 上，
\(h'(x)=1-\cos x\geqslant0\)，且 \(h(0)=0\)，所以 \(\sin x\leqslant x\).

再设 \(g(x)=\sin x-\frac{\sqrt2}{2}x\). 因为
\(g''(x)=-\sin x\leqslant0\)，所以 \(g\) 在 \([0,1]\) 上为凹函数，其最小值在端点取得. 又
\(g(0)=0\)，\(g(1)=\sin1-\frac{\sqrt2}{2}>0\)，
其中 \(1>\frac\pi4\) 且正弦函数在 \([0,\frac\pi2]\) 上递增. 因此 \(g(x)\geqslant0\)，即
\[
\frac{\sqrt2}{2}x\leqslant\sin x\leqslant x
\]

\item 记
\[
F_a(x)=ax+x^2+\frac{x^3}{2}+2(x+2)\cos x
\]
若题设不等式对所有 \(x\in[0,1]\) 成立，则对 \(x\in(0,1]\) 有
\[
a\leqslant \frac{4-x^2-x^3/2-2(x+2)\cos x}{x}
\]
将分子写成
\[
-2x+2(x+2)(1-\cos x)-x^2-\frac{x^3}{2}
\]
并令 \(x\to0^+\)，由 \(\frac{1-\cos x}{x}\to0\)，得到 \(a\leqslant-2\).

反之，若 \(a\leqslant-2\)，因 \(x\geqslant0\)，只需证明 \(F_{-2}(x)\leqslant4\). 而
\[
4-F_{-2}(x)=2(x+2)(1-\cos x)-x^2-\frac{x^3}{2}
\]
对 \(x/2\in[0,1]\) 使用第（1）问，得
\[
1-\cos x=2\sin^2\frac{x}{2}
 \geqslant2\left(\frac{\sqrt2}{2}\cdot\frac{x}{2}\right)^2
 =\frac{x^2}{4}
\]
因此
\[
4-F_{-2}(x)
 \geqslant2(x+2)\frac{x^2}{4}-x^2-\frac{x^3}{2}=0
\]
故实数 \(a\) 的取值范围为
\[
a\in(-\infty,-2]
\]
\end{enumerate}
\end{solution}

\section{选考题：解答题}

\begin{problem}
\begin{enumerate}
\item 如图，\(AB\) 为 \(\odot O\) 直径，直线 \(CD\) 与 \(\odot O\) 相切于 \(E\)，\(AD\) 垂直 \(CD\) 于 \(D\)，\(BC\) 垂直 \(CD\) 于 \(C\)，\(EF\) 垂直 \(AB\) 于 \(F\)，连接 \(AE\)，\(BE\). 证明：\(\angle FEB = \angle CEB\)；
\item 证明：\(EF^{2} = AD\cdot BC\).
\end{enumerate}
\bitmapfigure[float=inline,max width=0.20\linewidth]{img/2013/liaoning_liberal/q22_fig1.png}
\end{problem}

\begin{solution}
\begin{enumerate}
\item 因 \(AB\) 是直径，\(\angle AEB=90^\circ\). 又 \(EF\perp AB\)，所以在直角三角形
\(AEB\) 中，\(EF\) 是斜边 \(AB\) 上的高，得到
\[
\angle FEB=\angle EAB
\]
另一方面，切线弦定理给出
\[
\angle CEB=\angle EAB
\]
因此 \(\angle FEB=\angle CEB\).

\item 由 \(AD\perp CD\)，且 \(D\)，\(E\)，\(C\) 共线，三角形 \(ADE\) 在 \(D\) 点为直角. 由切线弦定理
\(\angle AED=\angle ABE\)，故
\[
AD=AE\sin\angle AED=AE\sin\angle ABE=AF
\]
最后一个等式来自直角三角形 \(AEB\) 中的三角函数关系.

同理，三角形 \(BCE\) 在 \(C\) 点为直角，且
\(\angle BEC=\angle BAE\)，所以
\[
BC=BE\sin\angle BEC=BE\sin\angle BAE=BF
\]
又由直角三角形 \(EAB\) 的斜边高定理，
\[
EF^2=AF\cdot FB
\]
结合 \(AD=AF\)、\(BC=BF\)，得
\[
EF^2=AD\cdot BC
\]
\end{enumerate}
\end{solution}

\begin{problem}
在直角坐标系 \(xOy\) 中，以 \(O\) 为极点，\(x\) 轴正半轴为极轴建立极坐标系. 圆 \(C_{1}\)，直线 \(C_{2}\) 的极坐标方程分别为 \(\rho = 4\sin \theta\)，\(\rho \cos \left(\theta - \frac{\pi}{4}\right) = 2\sqrt{2}\).

\begin{enumerate}
\item 求 \(C_{1}\) 与 \(C_{2}\) 交点的极坐标；
\item 设 \(P\) 为 \(C_{1}\) 的圆心，\(Q\) 为 \(C_{1}\) 与 \(C_{2}\) 交点连线的中点，已知直线 \(PQ\) 的参数方程为 \(\begin{cases}
x = t^3 + a,\\
y = \frac{b}{2} t^3 + 1,
\end{cases}\)（\(t \in \R\) 为参数），求 \(a\)，\(b\) 的值.
\end{enumerate}
\end{problem}

\begin{solution}
\begin{enumerate}
\item 由 \(\rho=4\sin\theta\)，得
\[
\rho^2=4\rho\sin\theta
\]
即 \(x^2+y^2=4y\).
由另一个方程，
\[
\rho\cos\left(\theta-\frac\pi4\right)=2\sqrt2
\]
化为 \(x+y=4\). 联立得
\[
x^2+(4-x)^2=4(4-x)
\]
即 \(2x(x-2)=0\). 所以交点为 \((0,4)\) 与 \((2,2)\)，对应的极坐标分别为
\(\left(4,\frac\pi2\right)\)，\(\left(2\sqrt2,\frac\pi4\right)\).

\item 圆 \(C_1:x^2+y^2=4y\) 的圆心为 \(P=(0,2)\). 两交点的中点为
\[
Q=\left(\frac{0+2}{2},\frac{4+2}{2}\right)=(1,3)
\]
因此直线 \(PQ\) 的方程为 \(y=x+2\). 由参数方程消去参数
\(t^3=x-a\)，得到
\[
y=\frac b2(x-a)+1
\]
与 \(y=x+2\) 比较斜率和截距，得
\[
\frac b2=1
\]
且 \(1-\frac{ab}{2}=2\)，
所以 \(b=2\)，\(a=-1\).
\end{enumerate}
\end{solution}

\begin{problem}
已知函数 \(f(x) = |x - a|\)，其中 \(a > 1\).

\begin{enumerate}
\item 当 \(a = 2\) 时，求不等式 \(f(x) \geqslant 4 - |x - 4|\) 的解集；
\item 已知关于 \(x\) 的不等式 \(\left|f(2x + a) - 2f(x)\right| \leqslant 2\) 的解集为 \(\{x \mid 1 \leqslant x \leqslant 2\}\)，求 \(a\) 的值.
\end{enumerate}
\end{problem}

\begin{solution}
\begin{enumerate}
\item 当 \(a=2\) 时，分三段讨论.
当 \(x\leqslant2\) 时，原不等式为
\[
2-x\geqslant x
\]
得 \(x\leqslant1\). 当 \(2\leqslant x\leqslant4\) 时，原不等式为
\(x-2\geqslant x\)，无解. 当 \(x\geqslant4\) 时，原不等式为
\[
x-2\geqslant8-x
\]
得 \(x\geqslant5\). 所以解集为
\[
(-\infty,1]\cup[5,+\infty)
\]

\item 因 \(a>1\)，有
\[
f(2x+a)=|2x|=2|x|
\]
题设不等式等价于
\[
\left||x|-|x-a|\right|\leqslant1
\]
当 \(x<0\) 时，\(|x|=-x\)，\(|x-a|=a-x\)，左边等于 \(a>1\)，无解.

当 \(0\leqslant x\leqslant a\) 时，左边为 \(|2x-a|\)，所以
\[
\frac{a-1}{2}\leqslant x\leqslant\frac{a+1}{2}
\]
当 \(x\geqslant a\) 时，左边等于 \(a>1\)，也无解. 因而解集为
\[
\left[\frac{a-1}{2},\frac{a+1}{2}\right]
\]
与题设的 \([1,2]\) 相同，故 \(a=3\).
\end{enumerate}
\end{solution}
